%%
%% spring-lecture05.tex
%% 
%% Made by Alex Nelson <pqnelson@gmail.com>
%% Login   <alex@lisp>
%% 
%% Started on  2026-04-09T12:13:15-0700
%% Last update 2026-04-09T12:13:15-0700
%% 

\lecture{}

\begin{definition}
For submodules $N\in S(M)$ and $P\in S(M)$, we say $P$ is a
\define{Complement} of $N$ if $N+P=M$ and $N\cap P=0$ (so $M=N\oplus P$).
\end{definition}

\begin{lemma}\label{lemma:2.14}
Let $M=\sum_{i\in J}N_{i}$ where each $N_{i}$ is a simple module. If
$P\in S(M)$, then there exists a subset $I\subset J$ such that
\begin{equation}
M=\left(\bigoplus_{i\in I}N_{i}\right)\oplus P.
\end{equation}
\end{lemma}

\begin{proof}
By Zorn's lemma: there exists $I\subset J$ such that $\{N_{i}\}_{i\in I}\cup\{P\}$
is maximal with respect to independence, i.e., $(\sum_{i\in I}N_{i})+P=(\bigoplus_{i\in I}N_{i})\oplus P$.
Let
\begin{equation}
M_{1}:=\left(\sum_{i\in I}N_{i}\right)+P.
\end{equation}
This maximality of $I$ implies $M_{1}\cap N_{j}\neq0$ for all $j\in J$.
But since the $N_{j}$ are simple, $N_{j}\subset M_{1}$ for all $j\in J$.
Then
\begin{equation}
M=\sum_{j\in J}N_{j}\subset M_{1}\subset M, 
\end{equation}
we see $M=M_{1}$.
\end{proof}

\begin{proposition}\label{prop:2.15}
Let $M$ be a right $A$-module. The following are equivalent:
\begin{enumerate}
\item $M$ is semisimple
\item $\displaystyle M=\sum_{\substack{N\in S(M)\\ N\text{ is simple}}}N$
\item $S(M)$ is a complemented lattice (i.e., for every $N\in S(M)$,
  $N$ has a complement in $S(M)$).
\end{enumerate}
\end{proposition}

\begin{proof}
$(1)\implies(2)$ clear

$(2)\implies(1)$ and $(2)\implies(3)$ by Lemma~\ref{lemma:2.14} (so
  $(1)\iff(2)$ and $(2)\implies (3)$).

$(3)\implies(2)$ First (3) is generalized to any submodule $P\in S(M)$
(i.e., $S(P)$ is a complemented lattice). If $M_{1}\in S(P)\subset S(M)$,
then there exists $M_{2}\in S(M)$ such that $M=M_{1}\oplus M_{2}$.
Then by the modular law
\begin{equation}
P=(M_{1}\oplus M_{2})\cap P=M_{1}\oplus(M_{2}\cap P),
\end{equation}
and since $M_{2}\cap P\in S(P)$, we see $S(P)$ is complemented.

Let $Q\in S(M)$ bebe a nontrivial proper submodule $Q\neq0$ and $Q\neq M$.
Let $u\in M\setminus Q$. Then by  Zorn's lemma, we may assume $Q$ is
maximal with respect to $u\notin Q$. Then by hypothesis (3),
$M=Q\oplus N$ where since $Q$ is nontrivial proper, we see $N\in S(M)$
is nontrivial proper.

Further, $u=w+v$ where $w\in Q$ and $v\in N$. If $N_{1}\in S(N)\setminus\{0,N\}$
is a nontrivial proper submodule of $N$, then $N=N_{1}\oplus N_{2}$
for some $N_{2}\neq0,N$. Then by maximality of $Q$, for only one of
$i=1$ or $i=2$ we have $u\in Q\oplus N_{i}$ and $u\in Q\oplus N_{i'}$ 
where $i'$ is the other possible value. [Pf: $M=Q\oplus N=Q\oplus(N_{1}\oplus N_{2})$
and $M=(Q\oplus N_{i})\oplus N_{i\pm 1}$.]
But it contradicts the maximality of $Q$ (i.e., $u\notin Q\oplus N_{i'}\propersupset Q$).
That is, $N$ is simple. Therefore by transfinite induction we obtain (2).
\end{proof}

\begin{corollary}\label{cor:2.16}
Let $M=\bigoplus_{i\in J}N_{i}$ be semisimple (so each $N_{i}$ is simple).
The following all hold:
\begin{enumerate}
\item For all $P\in S(M)$, both $P$ and $M/P$ are semisimple.
\item $S(M)$ is distributive if and only if for all $i\neq j$ we have
  $N_{i}\not\iso N_{j}$
\end{enumerate}
\end{corollary}

\begin{proof}
(1) By Proposition~\ref{prop:2.15}~(3) $M=P\oplus(M/P)$ and both $P$
  and $M/P$ are semisimple by (2).

(2) \forwardproof\ [Contrapositive] Observe for an $A$-module $N\neq0$
  and $Q=N\oplus N$, we have $S(Q)$ is not distributive. Let
  $N_{1}=N\oplus0\subset Q$ and $N_{2}=0\oplus N\subset Q$, and
  $N_{3}=\{(u,u)\in Q\mid u\in N\}\subset Q$. Then $N_{1}+N_{2}=Q$
  and $N_{1}\cap N_{3}=N_{2}\cap N_{3}=0$, so $N_{3}\iso N_{1}$
  and $N_{3}\cap(N_{1}+N_{2})=N_{3}$ but $(N_{3}\cap N_{1})+(N_{3}\cap N_{2})=0\neq N_{3}$
  so the distributive law does not hold. This observation proves the claim.

\backwardproof\ For any subset $I\subset J$, define
\begin{equation}
N(I) :=\sum_{i\in I}N_{i}.
\end{equation}
We will prove $\{N(I)\mid I\subset J\}$ is a distributive lattice.
Since $M=\bigoplus_{i\in J}N_{i}$, for any $I_{1}\subset I$ and
$I_{2}\subset I$ we have
\begin{equation}
N(I_{1}\cup I_{2})=N(I_{1})+N(I_{2})
\end{equation}
and
\begin{equation}
N(I_{1}\cap I_{2})=N(I_{1})\cap N(I_{2}).
\end{equation}
Then
\begin{subequations}
  \begin{align}
N(I_{1})\cap(N(I_{2})+N(I_{3}))
&= N\bigl(I_{1}\cap(I_{2}\cup I_{3})\bigr)\\
&= N\bigl((I_{1}\cap I_{2})\cup(I_{1}\cap I_{3})\bigr)\\
&=[N(I_{1})\cap N(I_{2})]+[N(I_{1})\cap N(I_{3})].
  \end{align}
\end{subequations}
Hence $\{N(I)\mid I\subset J\}$ is a distributive lattice.

We will show $S(M)=\{N(I)\mid I\subset J\}$. For any $P\in S(M)$, set
\begin{equation}
I=\{i\in J\mid P\cap N_{i}\neq\emptyset\}.
\end{equation}
In other words, $i\in I$ if and only if $N_{i}\subset P$;
equivalently, $N(I)\subset P$.

If
\begin{equation}\label{eq:math120c:spring2026:lecture5:cor2.16:big-if}
P\cap N(J\setminus I)=0,
\end{equation}
then by the modular law
\begin{subequations}
  \begin{align}
P &= (N(I)+N(J\setminus I))\cap P\\
&= N(I)+(N(J\setminus I)\cap P)\\
&= N(I) + 0 = N(I),
  \end{align}
\end{subequations}
and so
\begin{equation}
S(M)=\{N(I)\mid I\subset J\}.
\end{equation}
Now we just need to prove Equation~\eqref{eq:math120c:spring2026:lecture5:cor2.16:big-if}.

Assume that $P\cap N(J\setminus I)\neq0$ and choose the subset
$K\subset J\setminus I$ with the smallest cardinality such that $P\cap N(K)\neq0$.
Then you can easily check that $2\leq|K|<\infty$.
[Pf: Since $|K|=\min\{n\mid u=u_{i_{1}}+\dots+u_{i_{n}}\in P\cap\bigoplus^{n}_{j}N_{i_{j}},i_{j}\notin I\}<\infty$
but if $|K|=1$, then there exists $i\notin I$ such that $0\neq u=u_{i}\in P\cap N_{i}$
which contradicts $I=\{i\in J\mid P\cap N_{i}\neq\emptyset\}$.]
For $i\in K$, let $\pi_{i}\colon N(K)\to N_{i}$ be the
projection. Since
\begin{equation}
\ker(\pi|_{P\cap N(K)})\subset P\cap N(K\setminus\{i\}),
\end{equation}
and the minimality of $K$, we find that
\begin{equation}
\ker(\pi|_{P\cap N(K)})=0.
\end{equation}
That is to say $P\cap N(K)\iso N_{i}$. But this contradicts $|K|\geq2$
and $N_{i}\not\iso N_{j}$ for $i\neq j$.
\end{proof}

\begin{node}
Next, we will prove the uniqueness of the decomposition of semisimple
modules. This consists of two lemmas and one proposition.
\end{node}

\begin{lemma}\label{lemma:2.17}
Let $M=\bigoplus_{i\in J}N_{i}$ where each $N_{i}$ is simple and $N$
is simple such that there exists a nonzero $\phi\colon N\to M$.
Then there exists a $j\in J$ such that
$M\iso\phi(N)\oplus(\bigoplus_{i\neq j}N_{i})$ and $\phi(N)\iso N_{j}$.
\end{lemma}

\begin{proof}
By Lemma~\ref{lemma:2.14}, there exists a subset $J'\subset J$ such
that
\begin{equation}
M\iso\phi(N)\oplus\left(\bigoplus_{i\in J'}N_{i}\right).
\end{equation}
By Lemma~\ref{lemma:2.11},
\begin{subequations}
  \begin{align}
\bigoplus_{i\notin J'}N_{i} &\iso \frac{M}{\bigoplus_{i\in J'}N_{i}}\\
&\iso\phi(N)\\
&\iso N.
  \end{align}
\end{subequations}
Hence the claim.
\end{proof}

\begin{node}
  \begin{enumerate}
  \item In the next Lemma and Proposition, $\{N_{i}\mid i\in J\}$
is the set of representatives of simple $A$-modules.
\item For an $A$-module
$M$, we denote for $i\in J$,
\begin{equation}
M(i):=\sum_{\substack{N\in S(M)\\ N\iso N_{i}}}N.
\end{equation}
  \end{enumerate}
\end{node}

